\scanpage{74}
% archive foliation: 73
% lines 1--3
on devra expliciter dans le paradigme
le passage d'une description à
l'autre.
% lines 4--9
\deleted{La description}

1) Cas $J=\varnothing$\deleted{\inserted{$S\ne\varnothing$}}.
On se suppose, pour
notre description, que l'on n'est
pas dans les cas $g=0$, $\nu=0,1$
($g$ : genre de $\widehat X$, $S=\operatorname{Card}\nu$)
% lines 10--15
On associe à $(\widehat X,\widehat o,S)$\inserted{(ou) $(X,o)$} le
groupe fondamental $\Pi=\pi_1(X)$, considéré
comme groupe extérieur, mais
muni de sa $T$-structure :
locale indexée par $S$, i.e. une
\deleted{\unsure{famille}}\inserted{application}
% lines 16--17
\deleted{$\unsure{(\varphi_s)_{s\in S}|(\rho_s)_{s\in S}}$}
\margin{\textcircled{\unsure{1°}}}
\[S\longrightarrow\deleted{\unsure{g_s}}\operatorname{Homext}(T,\Pi)\]
% interlinear insertion following the local T-structure
\inserted{Ceci permet de définir $\Pi_1\ (\simeq\pi_1(\widehat X))$ comme
\unsure{groupe quot.} de $\Pi$ par le \unsure{ss-gpe inv.} engendré par les
$T\longrightarrow\Pi$.}
\inserted{Si $\nu=0$ i.e. $S=\varnothing$, il faut de}
% left margin, continuation of the preceding insertion
\inserted{plus à donner un iso
$T\longrightarrow\bigl[(\bigwedge^2\Pi_{1\mathrm{ab}})^\vee\bigr]^{\unsure{\mathrm{prim}|\mathrm{inv}}}$
de $T$ dans les formes
\unsure{primitives|invariantes} \unsure{ss} le groupe des
\unsure{autom.} ext. de \deleted{$\unsure{\Pi}$} $\Pi_1$.}
% lines 18--21
Ceci donne \inserted{$=$} une description
$\unsure{1}$-isotopique, et $=$ $\infty$-isotopique,
sauf si $g=0$, $\nu=2$ et $g=1$, $\nu=0$,
où les composantes \unsure{neutres}
